\chapter{Introduction}
\label{ch:intro}

The large-hadron collider (LHC) which is now in its third run, is as suggested by the name is 
an experiment which collides hadrons. This is precisely the scenario we look at in 
chapter~\ref{ch:hhZjet}. This thesis titled \emph{jet physics at the LHC and beyond} discusses
this process but also others, which leads us to the other topic we discuss which is the case of 
deep-inelastic scattering processes as in chapter~\ref{chapter:DIS}.

QCD like other QFTs varies in behaviour depending on scale. At low energies
the elementary particles bind together forming hadrons which are the final state particles that
we see in the detectors at the LHC and other experiments. At high energies these particles
behaves as free particles and so we can use the notion of perturbation theory, something that has
been widely used and produces remarkable results at high energies. In this thesis we focus on the high energy regime of QCD.

While perturbation theory offers us some nice results there are problems still in QCD. We investigate two
such problems in this thesis. The first is that within QCD (among other QFTs) objects known as jets are produced
these are messy objects which are formed by the clustering of many elementary QCD particles known as gluons. The notion of how these particles cluster together require some thought, which is brings up the topic
of jet algorithms. This is the idea of how a clustering procedure works such that we can mimic experiment in our theory. The second issue is large logarithms. These occur when the observable is small and can ruin perturbation theory.


\begin{enumerate}
	\item QCD complex theory due to the low energy physics
	\item jets are messy objects need to understand them at low energy
	\item need to understand clustering
	\item pQCD is great but there is the problem of resummation, need more results for scientific community
	
\end{enumerate}

This thesis discusses two topics within QCD. Firstly the notion of jet algorithms and secondly that of resummation.

In Chapter~\ref{ch:QCD}, we introduce the background context of the quantum chromodynamics. This will introduce the
relevant topics and key ideas crucial to understanding the original research presented in this thesis.
Chapter~\ref{ch:jet-phys} is a brief chapter explaining the concept of jets and how they must be dealt with such
leading us nicely into Chapter~\ref{chapter:DIS} presents a new jet algorithm used for deep-inelastic scattering processes.
In chapter~\ref{ch:resum}, we pivot to the resummation. In this chapter we introduce what this is exactly and provide the background 
relevant to understand the original research in chapter~\ref{ch:hhZjet}. Here, we present the \textsc{ARES} formalism, a semi-numerical resummation framework, for initial-state radiation for the very first time.
We conclude this thesis in chapter~\ref{ch:conc}, outlining how one could extend the research presented here
and how the findings contained have an impact on the scientific community.
